\documentclass[11pt,a4paper]{article}

\usepackage[T1]{fontenc}
\usepackage[utf8]{inputenc}
\usepackage{lmodern}
\usepackage[a4paper,margin=1in]{geometry}
\usepackage{xcolor}
\usepackage{hyperref}
\usepackage{graphicx}
\usepackage{booktabs}
\usepackage{longtable}
\usepackage{array}
\usepackage{enumitem}
\usepackage{titlesec}
\usepackage{tcolorbox}
\usepackage{listings}
\usepackage{setspace}

\definecolor{TitleBlue}{HTML}{163B65}
\definecolor{AccentOrange}{HTML}{C96A28}
\definecolor{SoftBlue}{HTML}{EAF2FB}
\definecolor{SoftGray}{HTML}{F4F6F8}
\definecolor{CodeBg}{HTML}{F7F7F7}
\definecolor{CodeFrame}{HTML}{D9D9D9}
\definecolor{LinkBlue}{HTML}{0B5CAD}

\hypersetup{
    colorlinks=true,
    linkcolor=TitleBlue,
    urlcolor=LinkBlue,
    citecolor=AccentOrange,
    pdftitle={SarfatiPelegChimera Windows Julia Guide},
    pdfauthor={Codex}
}

\lstdefinestyle{wincli}{
    backgroundcolor=\color{CodeBg},
    frame=single,
    rulecolor=\color{CodeFrame},
    basicstyle=\ttfamily\small,
    breaklines=true,
    columns=fullflexible,
    keepspaces=true,
    showstringspaces=false
}

\newtcolorbox{infobox}{
    colback=SoftBlue,
    colframe=TitleBlue,
    boxrule=0.8pt,
    arc=2pt
}

\newtcolorbox{warnbox}{
    colback=SoftGray,
    colframe=AccentOrange,
    boxrule=0.8pt,
    arc=2pt
}

\titleformat{\section}{\Large\bfseries\color{TitleBlue}}{\thesection}{0.75em}{}
\titleformat{\subsection}{\large\bfseries\color{AccentOrange}}{\thesubsection}{0.75em}{}

\setstretch{1.08}

\title{\textbf{\color{TitleBlue}SarfatiPelegChimera}\\[0.3em]
\large Julia Code Guide for Windows}
\author{Codex}
\date{August 2026}

\begin{document}
\maketitle
\tableofcontents
\bigskip

\section{Purpose of This Guide}
This document explains how to run the Julia-based codes in the
\texttt{SarfatiPelegChimera} package on a Windows machine. It is written for the
current Julia workflow only. It does not describe the Python, C++, or Mac-specific
parts of the broader Chimera workspace.

\begin{infobox}
The intended use of this package is to screen, validate, and compare delayed or
effective-coupling mechanisms that may generate two-group phase organization
resembling the Sarfati--Peleg-type phenomenology.
\end{infobox}

\section{What Is in the Package}
The package contains Julia scripts for:
\begin{itemize}[leftmargin=1.5em]
    \item coarse scans over parameter ranges,
    \item focused follow-up scans near promising regions,
    \item validation runs for top candidates,
    \item summary figure generation,
    \item CSV export of ranked parameter sets and diagnostics.
\end{itemize}

Typical outputs include:
\begin{itemize}[leftmargin=1.5em]
    \item \texttt{.csv} tables of scan results,
    \item summary \texttt{.png} figures,
    \item top-case diagnostic \texttt{.png} plots,
    \item validation summary \texttt{.png} plots.
\end{itemize}

\section{Windows Requirements}
You should have the following installed:
\begin{itemize}[leftmargin=1.5em]
    \item Julia 1.10 or newer,
    \item Visual Studio Code,
    \item the Julia extension for VS Code,
    \item a normal internet connection for the first package installation.
\end{itemize}

\subsection{Recommended Julia Packages}
Most scripts in this package rely on packages such as:
\begin{itemize}[leftmargin=1.5em]
    \item \texttt{DifferentialEquations}
    \item \texttt{CSV}
    \item \texttt{DataFrames}
    \item \texttt{Plots}
    \item \texttt{FFTW}
\end{itemize}

\section{Opening the Project in VS Code}
\begin{enumerate}[leftmargin=1.5em]
    \item Open VS Code.
    \item Choose \texttt{File > Open Folder}.
    \item Select the folder \texttt{SarfatiPelegChimera}.
    \item Open a terminal inside VS Code.
\end{enumerate}

\begin{warnbox}
All commands below should be run from the package folder itself, not from a parent
directory unless you explicitly pass the correct full path.
\end{warnbox}

\section{Preparing the Julia Environment}
Open a terminal in VS Code and run:

\begin{lstlisting}[style=wincli]
julia --project=.
\end{lstlisting}

Then, inside the Julia prompt, run:

\begin{lstlisting}[style=wincli]
using Pkg
Pkg.activate(".")
Pkg.instantiate()
\end{lstlisting}

If some packages are still missing, run:

\begin{lstlisting}[style=wincli]
using Pkg
Pkg.add("DifferentialEquations")
Pkg.add("CSV")
Pkg.add("DataFrames")
Pkg.add("Plots")
Pkg.add("FFTW")
\end{lstlisting}

\section{Running Julia Scripts}
To run a Julia script from the Windows terminal:

\begin{lstlisting}[style=wincli]
julia --project=. script_name.jl
\end{lstlisting}

Examples:

\begin{lstlisting}[style=wincli]
julia --project=. sarfati_peleg_stage2_julia.jl
julia --project=. sarfati_peleg_distributed_delay_search.jl
julia --project=. sarfati_peleg_pulse_coupled_search.jl
\end{lstlisting}

\subsection{Running a Validation Script}
If a script is intended to validate top candidates from an earlier CSV result,
run the screening script first and then run the validation script second.

\section{Using Multiple CPU Threads on Windows}
For faster scans, Julia can use multiple CPU threads.

\subsection{PowerShell}
\begin{lstlisting}[style=wincli]
$env:JULIA_NUM_THREADS="8"
julia --project=. script_name.jl
\end{lstlisting}

\subsection{Command Prompt}
\begin{lstlisting}[style=wincli]
set JULIA_NUM_THREADS=8
julia --project=. script_name.jl
\end{lstlisting}

\begin{infobox}
If your machine has more or fewer CPU cores, replace \texttt{8} with a suitable
number. In many cases, using all physical CPU cores gives a good speedup.
\end{infobox}

\section{Typical Workflow}
A practical workflow is:
\begin{enumerate}[leftmargin=1.5em]
    \item run a coarse scan,
    \item inspect the summary \texttt{.png} and top-ranked CSV entries,
    \item run the corresponding focused or validation script,
    \item inspect the top-case diagnostic plots,
    \item compare one-cluster and two-cluster indicators,
    \item refine the parameter window if necessary.
\end{enumerate}

\section{How to Read the Results}
Many of the scripts produce a mixture of:
\begin{itemize}[leftmargin=1.5em]
    \item score histograms,
    \item scatter plots of persistence versus antiphase error,
    \item heat maps of best score by parameter pair,
    \item detailed top-case diagnostics.
\end{itemize}

\subsection{Common Diagnostic Quantities}
Typical quantities used in the package include:
\begin{itemize}[leftmargin=1.5em]
    \item a one-cluster order parameter,
    \item a two-cluster order parameter,
    \item a group phase difference,
    \item a persistence score,
    \item an antiphase error,
    \item a final wrapped-phase snapshot,
    \item a late-time correlation matrix.
\end{itemize}

\section{Expected Output Files}
Depending on the script, you should expect outputs such as:
\begin{longtable}{>{\ttfamily}p{0.42\textwidth} p{0.5\textwidth}}
\toprule
File type & Meaning \\
\midrule
\textnormal{\texttt{*.csv}} & tabulated scan results and ranked parameter sets \\
\textnormal{\texttt{*_summary.png}} & overall summary plots for a scan \\
\textnormal{\texttt{*_top\_case\_NN.png}} & detailed diagnostics for high-ranking candidates \\
\textnormal{\texttt{*_top\_validation\_summary.png}} & summary of a validation stage \\
\bottomrule
\end{longtable}

\section{Troubleshooting}
\subsection{Package Not Found}
If Julia says a package is not found:
\begin{lstlisting}[style=wincli]
using Pkg
Pkg.activate(".")
Pkg.instantiate()
\end{lstlisting}

\subsection{Wrong Working Directory}
If scripts cannot find files, check that the terminal is opened in the
\texttt{SarfatiPelegChimera} folder before running them.

\subsection{Plots Not Appearing}
If plots are not displayed in the interactive window, they are often still saved
to disk as \texttt{.png} files. Check the folder contents after the run finishes.

\subsection{Long First Run}
The very first run may take noticeably longer because Julia is compiling packages.
Later runs are usually faster.

\section{Minimal Example Session}
\subsection{Step 1: Open Terminal in VS Code}
\begin{lstlisting}[style=wincli]
cd C:\path\to\SarfatiPelegChimera
\end{lstlisting}

\subsection{Step 2: Activate Environment}
\begin{lstlisting}[style=wincli]
julia --project=.
\end{lstlisting}

Then:
\begin{lstlisting}[style=wincli]
using Pkg
Pkg.activate(".")
Pkg.instantiate()
\end{lstlisting}

\subsection{Step 3: Run a Script}
\begin{lstlisting}[style=wincli]
julia --project=. sarfati_peleg_stage2_julia.jl
\end{lstlisting}

\subsection{Step 4: Inspect Output}
Look for generated CSV and PNG files in the package directory.

\section{Recommended Practice}
\begin{itemize}[leftmargin=1.5em]
    \item Keep the package folder structure unchanged.
    \item Run the coarse scan before any top-validation script.
    \item Keep notes of which script produced which result file.
    \item Use threads for broader scans.
    \item Re-run top candidates with a validation script rather than trusting the first screen alone.
\end{itemize}

\section{Final Note}
This guide is intended to be a stable Windows reference for the Julia part of the
\texttt{SarfatiPelegChimera} workflow. If the scripts evolve, the exact file names
may change, but the overall run pattern remains:
\[
\text{screen} \;\to\; \text{rank} \;\to\; \text{validate} \;\to\; \text{inspect}.
\]

\end{document}
